\begin{ouverture}[denombrement]{8}{Nombres premiers et petit théorème de Fermat}{Arithmétique}{Décomposition · Primalité · Crible · Chiffrement}
\objectif{Étudier la primalité et décomposer un entier.}
\objectif{Démontrer l’infinitude des nombres premiers.}
\objectif{Utiliser le petit théorème de Fermat.}
\objectif{Comprendre un chiffrement modulaire sur un exemple.}
\prerequis{PGCD, Bézout, Gauss et congruences.}
\end{ouverture}
\essentiel

\essSection{1}{Les briques multiplicatives des entiers}
Un nombre premier est un entier $p\geq2$ dont les seuls diviseurs positifs sont $1$ et $p$. Le nombre $1$ n'est pas premier.
Si $n\geq2$ est composé, $n=ab$ avec $2\leq a\leq b$ ; alors $a\leq\sqrt n$ et $a$ possède un diviseur premier. Pour tester la primalité, il suffit donc de chercher un diviseur premier au plus égal à $\sqrt n$.
\begin{theoreme}
Tout entier $n\geq2$ se décompose en produit de nombres premiers ; cette décomposition est unique à l'ordre des facteurs près.
\end{theoreme}
L'existence se prouve par récurrence forte : un entier est premier ou produit de deux entiers plus petits, que l'on décompose.
Pour l'unicité, Gauss implique qu'un premier divisant un produit divise un facteur. En comparant deux décompositions, on identifie puis simplifie un premier commun, et on poursuit.
\essSection{2}{Il y en a une infinité}
\begin{demonstration}[Preuve d'Euclide]
Supposer une liste finie $p_1,\ldots,p_r$ contenant tous les nombres premiers.
L'entier $N=p_1\cdots p_r+1>1$ possède un diviseur premier $q$.
Aucun $p_i$ ne divise $N$, car le reste serait $1$. Ainsi $q$ n'est pas dans la liste : contradiction.
On ne prétend pas que $N$ est toujours premier ; son diviseur premier suffit.
\end{demonstration}
\figurelivre{e08-crible.pdf}{Crible d'Ératosthène jusqu'à $50$ : garder chaque premier rencontré et éliminer ses multiples à partir de son carré. Les nombres conservés sont les nombres premiers.}
\begin{python}
def crible(n):
    premiers = [True]*(n+1)
    premiers[0] = False
    if n >= 1: premiers[1] = False
    p = 2
    while p*p <= n:
        if premiers[p]:
            for k in range(p*p, n+1, p):
                premiers[k] = False
        p += 1
    return [k for k in range(n+1) if premiers[k]]
\end{python}
On suppose $n\in\N$. Chaque composé possède un facteur premier au plus sa racine carrée, ce qui explique l'arrêt.
\essSection{3}{Fermat : une implication à respecter}
\begin{theoreme}[Petit théorème de Fermat]
Si $p$ est premier et $p\nmid a$, alors $a^{p-1}\equiv1\pmod p$.
Pour tout entier $a$, $a^p\equiv a\pmod p$.
\end{theoreme}
Une preuve possible : lorsque $p\nmid a$, les restes de $a,2a,\ldots,(p-1)a$ permutent les restes non nuls modulo $p$.
Le produit donne $a^{p-1}(p-1)!\equiv(p-1)!$ ; ce dernier facteur est inversible modulo $p$, donc on simplifie.
\begin{attention}
La réciproque est fausse. Si $a^{n-1}\not\equiv1\pmod n$ avec $a$ premier à $n$, on sait que $n$ est composé. Si la congruence vaut $1$, cela ne prouve pas que $n$ est premier.
Par exemple $341=11\times31$ et $2^{10}\equiv1\pmod{341}$, donc $2^{340}\equiv1$ bien que $341$ soit composé.
\end{attention}
\essSection{4}{Décomposer par un algorithme}
\begin{python}
def facteurs(n):
    resultat = []
    d = 2
    while d*d <= n:
        while n % d == 0:
            resultat.append(d)
            n //= d
        d += 1
    if n > 1: resultat.append(n)
    return resultat
\end{python}
On suppose $n\geq2$. Chaque facteur retiré est premier, car les plus petits facteurs possibles ont déjà été retirés. À la sortie, le dernier facteur éventuel est premier.
\essSection{5}{Un chiffrement affine, exemple pédagogique}
Coder les lettres de $0$ à $25$ ; chiffrer par $y\equiv ax+b\pmod{26}$.
Le déchiffrement est possible si $\operatorname{PGCD}(a,26)=1$ : avec un inverse $u$, $x\equiv u(y-b)\pmod{26}$.
Ces petits systèmes servent à comprendre l'arithmétique ; ils ne constituent pas une protection moderne. Le dossier RSA étudie un autre principe sur des nombres volontairement petits.

\exercices
\markright{Prendre confiance}
\begingroup\renewcommand{\otBesoinPlace}[1]{}\begin{minipage}{\linewidth}
\rubrique{Prendre confiance}
\begin{exercice}[1]{Premiers ou composés}

\begin{questions}
\item Classer $1,2,21,29$.
\item Pourquoi suffit-il de tester $2,3,5$ pour $29$ ?
\item Décomposer $84$.
\end{questions}
\end{exercice}
\begin{exercice}[1]{Décomposition et diviseurs}
$360=2^3\times3^2\times5$.
\begin{questions}
\item Combien a-t-il de diviseurs positifs ?
\item Donner le PGCD de $360$ et $84$.
\item Donner leur PPCM.
\end{questions}
\end{exercice}
\begin{exercice}[1]{Fermat direct}

\begin{questions}
\item Calculer $2^{12}$ modulo $13$ par Fermat.
\item En déduire $2^{100}$ modulo $13$.
\item Pourquoi $a=13$ ne permet-il pas d’utiliser la première forme du théorème modulo $13$ ?
\end{questions}
\end{exercice}
\begin{exercice}[1]{Un chiffrement}
$y\equiv5x+8\pmod{26}$, avec A codé $0$, B codé $1$.
\begin{questions}
\item Chiffrer A et B.
\item Vérifier que $21$ est l’inverse de $5$.
\item Écrire la formule de déchiffrement.
\end{questions}
\end{exercice}
\end{minipage}\endgroup
\clearpage
\markright{Appliquer et relier}
\begingroup\renewcommand{\otBesoinPlace}[1]{}\begin{minipage}{\linewidth}
\rubrique{Appliquer et relier}
\begin{exercice}[2]{Primalité de $223$}

\begin{questions}
\item Jusqu’à quelle borne faut-il chercher un diviseur premier ?
\item Tester les nombres premiers concernés.
\item Conclure.
\end{questions}
\end{exercice}
\begin{exercice}[2]{Crible}

\begin{questions}
\item Pourquoi éliminer les multiples de $p$ à partir de $p^2$ ?
\item Donner les nombres premiers au plus $20$.
\item Tester le programme pour $n=0$ et $n=1$.
\end{questions}
\end{exercice}
\begin{exercice}[2]{Décomposition algorithmique}

\begin{questions}
\item Appliquer \code{facteurs} à $180$.
\item Reconstituer le produit.
\item Pourquoi le dernier facteur résiduel, s’il est supérieur à $1$, est-il premier ?
\end{questions}
\end{exercice}
\begin{exercice}[2]{Un inverse par Fermat}

\begin{questions}
\item Utiliser Fermat pour exprimer un inverse de $a$ modulo un premier $p$ lorsque $p\nmid a$.
\item Calculer un inverse de $3$ modulo $7$.
\item Comparer à Bézout.
\end{questions}
\end{exercice}
\end{minipage}\endgroup
\clearpage
\markright{Chercher et argumenter}
\begingroup\renewcommand{\otBesoinPlace}[1]{}\begin{minipage}{\linewidth}
\rubrique{Chercher et argumenter}
\begin{exercice}[3]{Un faux test de primalité}
$341=11\times31$.
\begin{questions}
\item Verifier $2^{10}\equiv1\pmod{341}$.
\item En déduire $2^{340}\equiv1$.
\item Quelle conclusion erronée faut-il éviter ?
\end{questions}
\end{exercice}
\begin{exercice}[3]{L’infinité, sans faux raccourci}

\begin{questions}
\item Former $N=2\times3\times5\times7\times11\times13+1$.
\item Montrer qu’il peut être composé.
\item Pourquoi la preuve d’infinitude reste-t-elle valide ?
\end{questions}
\end{exercice}
\begin{exercice}[3]{RSA miniature, guidé}
On prend $p=5,q=11,n=55$ et $e=3$. Pour le message $m=7$, on chiffre $c\equiv m^3\pmod{55}$.
\begin{questions}
\item Calculer $c$.
\item Vérifier $3\times27\equiv1\pmod{40}$.
\item Calculer \code{pow(13, 27, 55)} et commenter.
\end{questions}
\end{exercice}
\begin{exercice}[3]{Racine rationnelle, prolongement}
$P(x)=a_nx^n+\cdots+a_0$ a des coefficients entiers. Une racine est $p/q$, fraction irréductible avec $q>0$.
\begin{questions}
\item Multiplier $P(p/q)=0$ par $q^n$.
\item Montrer que $q\mid a_n$.
\item Montrer que $p\mid a_0$ si $p\ne0$.
\end{questions}
\end{exercice}
\end{minipage}\endgroup
